\documentclass{article}
\usepackage[UTF8]{ctex}
\usepackage{amsmath,amssymb,amsfonts}
\usepackage{enumitem}
\usepackage{geometry}
\geometry{a4paper, margin=1in}
% === 错题本排版优化（不参与题目标识） ===
\setlength{\parindent}{0pt}
\setlength{\parskip}{0.35em plus 0.10em minus 0.08em}
\setlength{\emergencystretch}{3em}
\setlist{itemsep=0.25em, topsep=0.40em, parsep=0pt}
\allowdisplaybreaks[2]
% === 错题本排版优化结束 ===
\begin{document}
\section*{华南师范大学2026年高等代数真题}
\section*{151．华南师范大学2026年高等代数试卷}
一．填空题．每题5分，共30分．\\

\subsection{题目1}
1．设 2 是多项式 \(x^{4}-2 x^{3}+a x^{2}+b x-8\) 的二重根，求 \(a b=\)

\subsection{题目2}
2．设 \(A\) 是 \(n\) 阶方阵，\(B\) 是 \(m\) 阶方阵，记 \(D_{1}=\left|\begin{array}{ll}A & O \\ O & B\end{array}\right|, D_{2}=\left|\begin{array}{ll}O & A \\ B & O\end{array}\right|\) ，则 \(D_{1}: D_{2}=\)

\subsection{题目3}
3．实数域上秩为 \(n\) 的 \(n\) 元二次型 \(q\left(x_{1}, x_{2}, \cdots, x_{n}\right)\) 与 \(-q\left(x_{1}, x_{2}, \cdots, x_{n}\right)\) 等价，其中 \(n\) 是偶数，则\(q\left(x_{1}, x_{2}, \cdots, x_{n}\right)\) 的正惯性指数是 \(\_\_\_\_\)

\subsection{题目4}
4．设 2 阶可逆方阵 \(A\) 的特征多项式是 \(f(x)=x^{2}-10 x-24\) ，则 \(A^{-1}\) 的特征多项式是 \(\_\_\_\_\)

\subsection{题目5}
5．设 \(A=\left(\begin{array}{ccc}1 & 2 & -2 \\ 2 & 1 & 2 \\ 3 & 8 & 4\end{array}\right), b=\left(\begin{array}{l}x \\ 1 \\ 1\end{array}\right)\) ，若 \(A b\) 与 \(b\) 线性相关，则 \(x=\)

\subsection{题目6}
6．已知向量组

\[
\left\{\alpha_{1}, \alpha_{2}, \alpha_{3}\right\},\left\{\alpha_{1}, \alpha_{2}, \alpha_{3}, \alpha_{4}\right\},\left\{\alpha_{1}, \alpha_{2}, \alpha_{3}, \alpha_{5}\right\}
\]

的秩分别为3，3，4，则向量组 \(\left\{\alpha_{1}, \alpha_{2}, \alpha_{3}, \alpha_{4}-\alpha_{5}\right\}\) 的秩是 \(\_\_\_\_\)
\end{document}
